% Part: computability
% Chapter: computability-theory
% Section: halting-problem

\documentclass[../../../include/open-logic-section]{subfiles}

\begin{document}

\olfileid[fr]{cmp}{thy}{hlt}
\olsection{Le problème de l'arrêt}

Par construction, la fonction partielle calculable
universelle~$\fn{Un}(e,x)$ est définie si et
seulement si le calcul de la fonction de code~$e$
produit une valeur pour l'entrée~$x$. Il est
naturel de demander si l'on peut décider
effectivement quand cela se produit. La réponse
est négative. Dans le modèle des machines de
Turing, cela signifie que le problème de savoir
si une machine donnée s'arrête sur une entrée
donnée est indécidable par calcul. Le théorème
suivant exprime donc l'« indécidabilité du
problème de l'arrêt ». Nous en donnerons deux
preuves : la première prolonge la discussion
précédente, la seconde est plus directe.

\begin{thm}
\ollabel{thm:halting-problem}
Soit
\[
h(e, x)  =
\begin{cases}
1 & \text{si $\fn{Un}(e,x)$ est définie}\\
0 & \text{sinon.}
\end{cases}
\]
Alors $h$ n'est pas calculable.
\end{thm}

\begin{proof}
Supposons que $h$ soit calculable. Nous allons
voir que cela permettrait de définir une
fonction calculable totale universelle.
Posons
\[
\fn{Un'}(e,x) =
\begin{cases}
\fn{Un}(e,x) & \text{si $h(e,x)=1$}\\
0 & \text{sinon.}
\end{cases}
\]
La fonction $\fn{Un'}(e,x)$ serait alors totale,
et calculable si $h$ l'était. On peut par
exemple définir $g$ par récurrence primitive :
\begin{align*}
g(0, e, x) & \simeq 0 \\
g(y+1, e, x) & \simeq \fn{Un}(e,x);
\end{align*}
Cette récurrence conserve $g(0,e,x)=0$ même si $\fn{Un}(e,x)$ diverge ; elle ne calcule $\fn{Un}(e,x)$ qu'aux rangs positifs. On obtient alors
\[
\fn{Un'}(e,x) \simeq g(h(e,x),e,x).
\]
Comme $\fn{Un'}(e,x)$ coïncide avec
$\fn{Un}(e,x)$ là où cette dernière est
définie, $\fn{Un'}$ serait universelle
pour les fonctions partielles calculables
qui se trouvent être totales. Cela contredit le
\olref[nou]{thm:no-univ}.
\end{proof}

\begin{proof}
Supposons encore $h(e,x)$ calculable.
Définissons la fonction~$g$ par
\[
g(x) =
\begin{cases}
  0           & \text{si $h(x,x)=0$}\\
  \fundefined & \text{sinon.}
\end{cases}
\]
La fonction $g$ est partielle calculable.
On peut, par exemple, la définir par
$\umin{y}{h(x,x)=0}$. Elle possède donc
un indice~$e$ tel que
$g(x) \simeq \fn{Un}(e,x)$ pour tout~$x$.
La valeur $g(e)$ est-elle définie ?
Si elle l'est, la définition de~$g$
donne $h(e,e)=0$ (puisque $h$ est totale
par définition et ne prend que les
valeurs~$0$ et~$1$). Par définition de~$h$,
$\fn{Un}(e,e)$ est alors indéfinie.
L'égalité extensionnelle
$g(x) \simeq \fn{Un}(e,x)$ implique que
$g(e)$ est indéfinie : contradiction.
Si, au contraire, $g(e)$ est indéfinie,
alors $h(e,e)\neq 0$, donc $h(e,e)=1$.
Il s'ensuit que $\fn{Un}(e,e)$ est définie
et, par la même égalité extensionnelle,
que $g(e)$ est définie : nouvelle
contradiction.
\end{proof}

\begin{tagblock}{TMs}
\begin{explain}
On peut aussi exprimer cet argument à l'aide
de machines de Turing. Supposons qu'une machine~$H$
reçoive la description d'une machine~$E$ et
une entrée~$x$, et décide si $E$ s'arrête
sur~$x$. Nous pourrions alors construire
une machine~$G$ qui, recevant une seule
entrée~$x$, fait exécuter~$H$ afin de savoir
si la machine $M_x$ d'indice~$x$ s'arrête
sur l'entrée~$x$, puis adopte le comportement
opposé. Si $H$ annonce que $M_x$ s'arrête,
$G$ entre dans une boucle infinie; si
$H$ annonce que $M_x$ ne s'arrête pas,
$G$ s'arrête aussitôt. Que ferait $G$
sur son propre indice ? Elle s'arrêterait
si et seulement si elle ne s'arrêtait
pas : contradiction. Une telle machine~$H$
n'existe donc pas.
\end{explain}
\end{tagblock}

\end{document}
